%============================================================
\chapter{Follow-the-Gap-Algorithmus}
\label{chap:algorithmus}
%============================================================

\section{Grundidee}
\label{sec:ftg_idea}

Der \emph{Follow-the-Gap}-Algorithmus (FTG) ist ein reaktives
Verfahren zur Hindernissvermeidung, das ausschließlich auf den
Rohdaten eines 2D-LiDAR-Scanners arbeitet.
Die Grundidee besteht darin, im aktuellen LiDAR-Scan die größte
freie Lücke (\enquote{Gap}) zu identifizieren und das Fahrzeug
auf deren Mittelpunkt auszurichten.

\begin{enumerate}[label=\arabic*., leftmargin=*]
  \item \textbf{Sicherheitsblase erzeugen:}
        Um jeden detektierten Messpunkt, der näher als ein
        Schwellwert \(r_{\min}\) liegt, wird eine virtuelle
        Sicherheitsblase mit Radius \(r_{\text{bubble}}\) gelegt.
        Alle Messpunkte innerhalb dieser Blase werden auf
        \(0\) gesetzt.
  \item \textbf{Freie Lücken bestimmen:}
        Zusammenhängende Folgen von Messpunkten mit Abstand
        \(> r_{\min}\) bilden sogenannte \emph{freie Lücken}.
  \item \textbf{Beste Lücke wählen:}
        Die breiteste Lücke (größte Anzahl aufeinanderfolgender
        freier Messpunkte) wird als Ziellücke ausgewählt.
  \item \textbf{Lenkwinkel berechnen:}
        Der Lenkwinkel \(\delta\) wird so gewählt, dass das Fahrzeug
        auf den \emph{Mittelpunkt} der besten Lücke zeigt.
  \item \textbf{Geschwindigkeitsregelung:}
        Die Fahrgeschwindigkeit \(v\) wird in Abhängigkeit des
        berechneten Lenkwinkels reduziert, um bei engen Kurven
        stabiler zu fahren.
\end{enumerate}

%------------------------------------------------------------
\section{Mathematische Beschreibung}
\label{sec:ftg_math}
%------------------------------------------------------------

Sei \(\mathcal{S} = \{s_0, s_1, \ldots, s_{N-1}\}\) die Folge der
\(N\) LiDAR-Messpunkte eines einzelnen Scans, wobei
\(s_i = (r_i, \theta_i)\) den Abstand \(r_i\) und den Winkel
\(\theta_i\) des \(i\)-ten Strahls bezeichnet.

\subsection{Sicherheitsblase}
\label{subsec:bubble}

Sei \(i^* = \arg\min_i r_i\) der Index des nächsten Hindernispunkts.
Alle Messpunkte \(s_j\) mit
\[
  \left| \theta_j - \theta_{i^*} \right| \leq
  \arctan\!\left(\frac{r_{\text{bubble}}}{r_{i^*}}\right)
\]
werden auf \(r_j \leftarrow 0\) gesetzt.
Der Parameter \(r_{\text{bubble}}\) entspricht etwa der halben
Fahrzeugbreite zuzüglich eines Sicherheitszuschlags.

\subsection{Spaltenerkennung}
\label{subsec:gaps}

Eine freie Lücke \(G_k = [i_{\text{start}}^{(k)},\,
i_{\text{end}}^{(k)}]\) ist eine maximale Folge von Indizes, für
die gilt:
\[
  r_i > r_{\min}
  \quad \forall\, i \in
  \bigl[i_{\text{start}}^{(k)},\, i_{\text{end}}^{(k)}\bigr].
\]

Die \emph{beste Lücke} \(G^*\) wird durch
\[
  G^* = \arg\max_{G_k}
  \bigl(i_{\text{end}}^{(k)} - i_{\text{start}}^{(k)}\bigr)
\]
bestimmt.

\subsection{Lenkwinkelberechnung}
\label{subsec:steering}

Der Zielpunkt liegt im Mittelpunkt der besten Lücke:
\[
  i_{\text{goal}} =
  \left\lfloor
    \frac{i_{\text{start}}^* + i_{\text{end}}^*}{2}
  \right\rfloor.
\]

Der entsprechende Lenkwinkel ergibt sich direkt aus dem Winkel
des Zielpunkts:
\[
  \delta = \theta_{i_{\text{goal}}}.
\]

\subsection{Geschwindigkeitsprofil}
\label{subsec:velocity}

Um eine stabile Kurvenfahrt zu gewährleisten, wird die Zielgeschwindigkeit
\(v\) als nichtlineare Funktion des Lenkwinkels gewählt:
\[
  v(\delta) =
  \begin{cases}
    v_{\max}          & \text{falls } |\delta| \leq \alpha_1, \\
    v_{\text{mid}}    & \text{falls } \alpha_1 < |\delta| \leq \alpha_2, \\
    v_{\min}          & \text{falls } |\delta| > \alpha_2,
  \end{cases}
\]
wobei \(\alpha_1 = 10°\), \(\alpha_2 = 20°\) in unserer Implementierung
verwendet wurden (vgl.\ \cref{tab:parameter}).

%------------------------------------------------------------
\section{Diskussion der Strategie}
\label{sec:ftg_discussion}
%------------------------------------------------------------

\textbf{Vorteile:}
\begin{itemize}
  \item Geringer Rechenaufwand – echtzeitfähig bei $>40\,\text{Hz}$.
  \item Keine Karte oder Vorwissen über die Umgebung notwendig.
  \item Robust gegenüber Messrauschen dank Radiusmittelung.
\end{itemize}

\textbf{Nachteile und bekannte Grenzen:}
\begin{itemize}
  \item Rein reaktiv -- keine Vorausplanung möglich.
  \item Kann in engen Korridoren oder bei konkaven Hindernisformen
        in lokalen Minima steckenbleiben.
  \item Die Wahl der Parameter (\(r_{\min}\), \(r_{\text{bubble}}\),
        Schwellwerte für \(v\)) ist stark umgebungsabhängig und
        erfordert aufwändiges Tuning.
  \item In unserem Aufbau führte die begrenzte Scan-Sichtweite in
        engen Kurven zu Fehlklassifizierungen freier Lücken
        (siehe \cref{chap:ergebnisse}).
\end{itemize}
